接着考虑原函数的存在性. 当然, 不能定义广义函数的定积分. 尽管如此, 仍能证明原函数存在.

\begin{Theorem}
\label{thm:distribution-primitive-existence}
设 $I=(a,b)$ 为 $\mathbb{R}$ 中的开区间, 且 $f\in\mathcal{D}^{\prime}(I)$. 那么存在 $u\in\mathcal{D}^{\prime}(I)$, 使 $u'=f$. 原函数 $u$ 在相差一个常数的意义下唯一.
\end{Theorem}

\begin{Proof}
唯一性部分由前一定理显然得出. 为构造原函数, 使用分解 \eqref{eq:test-function-one-dimensional-zero-integral-decomposition}
\begin{equation}\label{eq:primitive-construction-test-decomposition}
\phi(x)=\phi_1(x)\int_I\phi(s)\,\mathrm{d}s+\psi(x),
\end{equation}
并令
\begin{equation}\label{eq:primitive-construction-zero-integral-test-primitive}
\chi(x)=\int_a^x\psi(y)\,\mathrm{d}y.
\end{equation}
然后定义
\begin{equation}\label{eq:primitive-distribution-definition-pairing}
(u,\phi)=C\int_I\phi(x)\,\mathrm{d}x-(f,\chi),
\end{equation}
其中 $C$ 为任意常数. 若 $\phi=\eta'$, 则 $\int_I\phi(x)\,\mathrm{d}x=0$ 且 $\psi=\phi$; 因而 $\chi=\eta$. 由此得到
\begin{equation}\label{eq:primitive-distribution-derivative-check}
(u,\eta')=-(f,\eta);
\end{equation}
因而 $u'=f$.
\end{Proof}

单连通区域上任意无旋向量场都是梯度这一多维结果也可推广到分布; 其证明远比一维情形复杂, 这里不予给出.

求解常微分方程最基本的技巧, 是将其化为 $y'=f$ 的形式; 这就是为什么求解常微分方程被称为对它作"积分". 这些步骤对分布解也适用. 例如考虑常微分方程
\begin{equation}\label{eq:distribution-linear-first-order-ode}
y'=a(x)y+f(x).
\end{equation}


假设 $a\in C^\infty(\mathbb{R})$ 且 $f\in\mathcal{D}^{\prime}(\mathbb{R})$. 现在可以令
\begin{equation}\label{eq:distribution-ode-integrating-factor-substitution}
y(x)=z(x)\exp\left(\int_0^xa(s)\,\mathrm{d}s\right);
\end{equation}
注意, 分布乘以 $C^\infty$ 函数有良好定义, 并且容易证明微分的乘积法则成立. 于是得到新的常微分方程
\begin{equation}\label{eq:distribution-ode-primitive-reduction}
z'(x)=f(x)\exp\left(-\int_0^xa(s)\,\mathrm{d}s\right).
\end{equation}


由 Theorem~\ref{thm:distribution-primitive-existence}, 知道这一常微分方程具有一个单参数解族.

特别地, 若 $f$ 为连续函数, 则 \eqref{eq:distribution-linear-first-order-ode} 的所有分布解都是经典解. 对奇异常微分方程, 这不一定成立; 例如常数 1 和 Heaviside 函数都是 $xy'=0$ 的解.
